\section{Model, profile target and penalised estimation}\label{sec:model}

\subsection{Clustered quantile model and asymptotic regime}\label{subsec:model}

Let $D_i=\{(Y_{ij},X_{ij},Z_{ij}):j=1,\ldots,m_i\}$ denote cluster $i$, where $i=1,\ldots,n$. The clusters are independent, while observations within a cluster may be arbitrarily dependent subject to the moment conditions stated later. Here $Y_{ij}\in\R$ is the response, $X_{ij}\in\R^p$ is the fixed-effect design, and $Z_{ij}\in\R^q$ is the design for cluster-specific effects. We consider
\begin{equation}\label{eq:structural-model}
 Q_{\tau}(Y_{ij}\mid X_i,Z_i,b_i)
 =X_{ij}^{\mathsf T}\beta^0+Z_{ij}^{\mathsf T}b_i,
 \qquad \tau\in(0,1),
\end{equation}
where $X_i=(X_{i1},\ldots,X_{im_i})^{\mathsf T}$ and $Z_i=(Z_{i1},\ldots,Z_{im_i})^{\mathsf T}$. The dimension $q$ is fixed. The number of clusters $n$ tends to infinity, the cluster sizes satisfy $m_i\leq m_{\max}<\infty$, and $p=p_n$ may exceed $n$. The structural representation in \eqref{eq:structural-model} motivates the working criterion, but the primary inferential target is the regularised profile parameter defined below.

We use the number of independent clusters, rather than the total number of observations, as the asymptotic sample size. Each cluster receives equal weight by averaging its within-cluster loss. This convention avoids changing the population target when cluster sizes vary. Observation-weighted versions are obtained by replacing $m_i^{-1}$ by prespecified bounded weights; their profile calculus is identical after carrying the weights through the definitions.

The fixed-effect vector may include a fixed number of always-included coefficients, such as an intercept, time, treatment, or demographic covariates. In computation these coordinates receive penalty weight zero. To keep notation concise, the theoretical statements write a single vector $\beta$ and impose sparsity on its high-dimensional component.

\subsection{An unsmoothed regularised profile target}\label{subsec:target}

For a positive-definite $q\times q$ working penalty matrix $\Lambda$, define the cluster criterion
\begin{equation}\label{eq:cluster-unsmoothed}
 \ell_i(\beta,\gamma;D_i)
 =\frac{1}{m_i}\sum_{j=1}^{m_i}
 \rho_{\tau}\!\left(Y_{ij}-X_{ij}^{\mathsf T}\beta-Z_{ij}^{\mathsf T}\gamma\right)
 +\frac{1}{2}\gamma^{\mathsf T}\Lambda\gamma,
\end{equation}
where $\rho_{\tau}(u)=u\{\tau-I(u<0)\}$. Since $\Lambda$ is positive definite, the nuisance minimiser
\begin{equation}\label{eq:gamma-target}
 \gamma_i^{\dagger}(\beta;D_i)
 =\argmin_{\gamma\in\R^q}\ell_i(\beta,\gamma;D_i)
\end{equation}
exists and is unique for every $\beta$ and realised cluster. The population profile risk is
\begin{equation}\label{eq:population-profile-risk}
 R(\beta)=\E\left[\ell_i\{\beta,\gamma_i^{\dagger}(\beta;D_i);D_i\}\right],
 \qquad
 \beta^{\star}=\argmin_{\beta\in\R^p}R(\beta).
\end{equation}
The expectation is over an independent generic cluster. The parameter $\beta^{\star}$ depends on $\Lambda$ but not on a smoothing bandwidth. We call it the regularised profile parameter.

The distinction between $\beta^{\star}$ and $\beta^0$ is substantive. When cluster sizes remain bounded, the data contain limited information about each $b_i$. The ridge subproblem replaces the unidentified or unstable cluster effect by a regularised cluster projection, and $\beta^{\star}$ records the fixed-effect component of the resulting population projection. Equality $\beta^{\star}=\beta^0$ requires an orthogonality condition,
\begin{equation}\label{eq:structural-orthogonality}
 \E\left[\partial_{\beta}
 \ell_i\{\beta^0,\gamma_i^{\dagger}(\beta^0;D_i);D_i\}\right]=0,
\end{equation}
interpreted through a subgradient when needed. This condition can hold for within-cluster centred slope covariates in a random-intercept location-shift model, but it does not generally hold for the intercept, between-cluster covariates, random-slope misspecification, or arbitrary shrinkage. We therefore state all inferential results for $\beta^{\star}$ and discuss structural interpretation coordinate by coordinate.

The matrix $\Lambda$ may be written as $\lambda_{\gamma}G^{-1}$, where $G$ is a prespecified positive-definite working correlation matrix for the cluster effects. The main theory treats $\Lambda$ as fixed with bounded eigenvalues. Estimating a low-dimensional $G$ is a practically useful extension, but its uncertainty must be included in the profile linearisation before it is used in the main inferential theorem.

\subsection{Convolution smoothing and the penalised estimator}\label{subsec:estimation}

Let $K$ be a symmetric second-order probability density supported on a compact interval, and set $K_h(u)=h^{-1}K(u/h)$. The convolution-smoothed check loss is
\begin{equation}\label{eq:smoothed-loss}
 \rho_{\tau,h}(u)=\int \rho_{\tau}(u-hv)K(v)\,dv.
\end{equation}
Write $\psi_{\tau,h}(u)=\rho_{\tau,h}'(u)$ and $\phi_{\tau,h}(u)=\rho_{\tau,h}''(u)=K_h(u)$. The smoothed cluster criterion is
\begin{equation}\label{eq:cluster-smoothed}
 \ell_{i,h}(\beta,\gamma;D_i)
 =\frac{1}{m_i}\sum_{j=1}^{m_i}
 \rho_{\tau,h}\!\left(Y_{ij}-X_{ij}^{\mathsf T}\beta-Z_{ij}^{\mathsf T}\gamma\right)
 +\frac{1}{2}\gamma^{\mathsf T}\Lambda\gamma.
\end{equation}
Its population profile minimiser is denoted by $\beta_h^{\star}$. The estimator is
\begin{equation}\label{eq:penalised-estimator}
 (\widehat\beta,\widehat\gamma)
 \in\argmin_{\beta\in\R^p,\,\gamma_i\in\R^q}
 \left[
 \frac{1}{n}\sum_{i=1}^n\ell_{i,h}(\beta,\gamma_i;D_i)
 +\lambda_{\beta}\sum_{k=1}^p w_k|\beta_k|
 \right],
\end{equation}
where $w_k=0$ for always-included coordinates and $w_k=1$ for penalised coordinates. The criterion is jointly convex. The $\gamma_i$ subproblems are strictly convex and separable across clusters.

A block algorithm alternates between a proximal-gradient update for $\beta$ and low-dimensional Newton updates for the $\gamma_i$. Because the objective is convex, standard sufficient-decrease and vanishing-inner-error conditions are enough for convergence to a global minimiser \cite{tseng2009,razaviyayn2013,wright2015}. We do not require global strong convexity of the fixed-effect block. Algorithm~\ref{alg:estimation} records the implementation used for the corrected numerical study.

\begin{algorithm}[t]
\caption{Smoothed regularised profile estimation}\label{alg:estimation}
\begin{algorithmic}[1]
\State Initialise $\beta^{(0)}$ by pooled smoothed quantile Lasso and set $\gamma_i^{(0)}=0$.
\For{$t=0,1,\ldots$}
  \State Update $\beta$ by proximal gradient/FISTA for the criterion in \eqref{eq:penalised-estimator} with $\gamma=\gamma^{(t)}$.
  \For{$i=1,\ldots,n$ in parallel}
    \State Update $\gamma_i$ by damped Newton minimisation of $\ell_{i,h}(\beta^{(t+1)},\gamma_i;D_i)$.
  \EndFor
  \State Stop when the relative objective change and the maximum KKT residual are below tolerance.
\EndFor
\State Return $(\widehat\beta,\widehat\gamma)$ and retain the complete optimisation log.
\end{algorithmic}
\end{algorithm}

The theoretical order $\lambda_{\beta}\asymp\sqrt{\log(p)/n}$ supplies a reference grid. In finite samples, $\lambda_{\beta}$ and a scalar $\lambda_{\gamma}$ are selected by subject-level cross-validation, never by splitting individual observations within a cluster. The bandwidth is parameterised as $h=c_h\widehat s_Yn^{-1/3}$, where $\widehat s_Y$ is a robust response scale estimated within the training clusters and $c_h$ is varied over a fixed sensitivity grid. This order is compatible with the explicit feasible regime in Section~\ref{subsec:feasible}.

\subsection{Exact profile score and effective Hessian}\label{subsec:profile-calculus}

For each $\beta$, let
\begin{equation}\label{eq:empirical-gamma-profile}
 \widehat\gamma_i(\beta)
 =\argmin_{\gamma\in\R^q}\ell_{i,h}(\beta,\gamma;D_i)
\end{equation}
and define
\begin{equation}\label{eq:empirical-profile-loss}
 L_{n,h}^{\mathrm{prof}}(\beta)
 =\frac{1}{n}\sum_{i=1}^n
 \ell_{i,h}\{\beta,\widehat\gamma_i(\beta);D_i\}.
\end{equation}
Let $r_i(\beta)=Y_i-X_i\beta-Z_i\widehat\gamma_i(\beta)$,
\begin{equation}\label{eq:WiBiAi}
 W_i(\beta)=\frac{1}{m_i}\diag\{\phi_{\tau,h}(r_{ij}(\beta))\}_{j=1}^{m_i},
 \quad
 B_i(\beta)=Z_i^{\mathsf T}W_i(\beta)Z_i+\Lambda,
\end{equation}
\begin{equation}\label{eq:Ai}
 A_i(\beta)=B_i(\beta)^{-1}Z_i^{\mathsf T}W_i(\beta)X_i,
 \qquad
 \widetilde X_i(\beta)=X_i-Z_iA_i(\beta).
\end{equation}

\begin{proposition}[Exact profile calculus]\label{prop:profile-calculus}
For every $\beta$ at which the smoothed criterion is twice differentiable, the gradient of the complete ridge-profile loss is
\begin{equation}\label{eq:profile-score}
 g_{n,h}(\beta)=\nabla L_{n,h}^{\mathrm{prof}}(\beta)
 =-\frac{1}{n}\sum_{i=1}^n\frac{1}{m_i}
 X_i^{\mathsf T}\psi_{\tau,h}\{r_i(\beta)\}.
\end{equation}
Its Hessian is
\begin{align}
 \widehat H_{\mathrm{eff}}(\beta)
 &=\frac{1}{n}\sum_{i=1}^n
 \left[
 X_i^{\mathsf T}W_iX_i
 -X_i^{\mathsf T}W_iZ_iB_i^{-1}Z_i^{\mathsf T}W_iX_i
 \right] \label{eq:heff-schur}\\
 &=\frac{1}{n}\sum_{i=1}^n
 \left[
 \widetilde X_i^{\mathsf T}W_i\widetilde X_i
 +A_i^{\mathsf T}\Lambda A_i
 \right]. \label{eq:heff-ridge-identity}
\end{align}
All matrices on the right-hand sides are evaluated at $\beta$. The two expressions are identical.
\end{proposition}

The proof follows from the envelope theorem, implicit differentiation of the $\gamma_i$ stationarity equations, and direct algebra; it is given in the \supp. Proposition~\ref{prop:profile-calculus} fixes the object used throughout the paper. The score in \eqref{eq:profile-score} uses $X_i$, not $\widetilde X_i$. Residualisation enters the Hessian through the derivative of the profile map. When the Schur complement is written as a weighted Gram matrix of $\widetilde X_i$, the positive term $A_i^{\mathsf T}\Lambda A_i$ is indispensable. It vanishes only when the nuisance penalty vanishes or the fixed and nuisance designs are locally orthogonal.

The implementation evaluates three identities before any simulation or data analysis is accepted: a finite-difference check of \eqref{eq:profile-score}, a finite-difference check of \eqref{eq:heff-schur}, and a direct numerical comparison of \eqref{eq:heff-schur} with \eqref{eq:heff-ridge-identity}. These checks isolate algebraic errors from asymptotic approximation errors.

